\appendix

\section{Proofs}\label{app:proofs}

The identities, monotonicity statements, and optimality inequalities appearing in these proofs---together with the application arithmetic of Section~\ref{sec:gaza}, over exact rationals---have additionally been machine-checked in Lean~4 with mathlib; the development compiles with no unproven statements and is in the \texttt{formal/} directory of the companion repository, which documents exactly which statement corresponds to which theorem below. Steps involving posteriors, quantiles, and limits are formalised only in simplified kernel form there; for those, the proofs below are the complete argument, and they are additionally exercised by the symbolic and Monte-Carlo checks in the repository's verification scripts.

\begin{proof}[Proof of Theorem~\ref{thm:bounds}]
Throughout, $w>0$ and $f\le a$, so $w+\mu(a-f)>0$ for every $\mu\ge 0$ and all denominators below are positive; the inputs $(\omega,w,a,f)$ are treated as population quantities (sampling error enters separately through the delta method of Section~\ref{sec:bounds}). Let $\lambda$ be the civilian per-capita death risk over the analysis window in the women-and-children class $W$ (a population-level proportion, not an instantaneous rate; deaths are treated as proportional to class populations and integer rounding is ignored), so class $\mathrm{AM}$ civilians have risk $\mu\lambda$. Expected civilian deaths in $W$ and in civilian-$\mathrm{AM}$ are $\lambda w N$ and $\lambda\mu(a-f)N$, so the share of \emph{civilian} deaths landing in $W$ is $w/(w+\mu(a-f))$. Combatant deaths land in $\mathrm{AM}$, so only the civilian fraction $1-q$ of deaths reaches $W$:
\[
\omega=(1-q)\,\frac{w}{w+\mu(a-f)},
\]
which solves to \eqref{eq:qmu}; moreover $\partial_\mu q=-\omega(a-f)/w\le 0$, so $q(\mu)$ is nonincreasing (strictly decreasing unless $\omega=0$ or $a=f$).

\emph{Attainability.} Fix $\mu^\ast\in[\mu_{\mathrm{lo}},\mu_{\mathrm{hi}}]$ with $q(\mu^\ast)\in(0,1)$. With $(D,N)$ observed and fixed, the $W$-class risk is pinned at $\lambda=\omega D/(wN)$, not freely selectable; the feasibility conditions of Theorem~\ref{thm:bounds} ($\omega D\le wN$ and $\mu_{\mathrm{hi}}\,\omega D\le wN$) guarantee $\lambda\le 1$ and $\mu^\ast\lambda\le 1$, so no class's implied deaths exceed its population. The configuration with combatant share $q(\mu^\ast)$ and risks $(\lambda,\mu^\ast\lambda)$ is then a valid allocation and reproduces the observed $\omega$ exactly, so $q(\mu^\ast)$ cannot be excluded by the data. Boundary values are attained in the closure: $q(\mu^\ast)=1$ forces $\omega=0$ (no civilian deaths, the limit $\lambda\downarrow 0$), and $q(\mu^\ast)=0$ is attained directly with $\lambda>0$ and no combatant deaths.

\emph{Sharpness.} Conversely, any configuration consistent with the model has some $\mu\in[\mu_{\mathrm{lo}},\mu_{\mathrm{hi}}]$, and the display then forces $q=q(\mu)$. Since $\mu\mapsto q(\mu)$ is continuous and monotone, the set of attainable values is exactly the interval $[q(\mu_{\mathrm{hi}}),q(\mu_{\mathrm{lo}})]$; intersecting with $[0,1]$, the logical range of a share, gives the identified set. The benchmark identity \eqref{eq:ident} is the case $\mu=1$, using $w+a=1$. Corollary~\ref{cor:popbudget} follows since $D_M\le M$ and $D>0$ give $q=D_M/D\le M/D$; the joint set is the intersection of two intervals and is empty exactly when the stated endpoints fail to be ordered.
\end{proof}

\begin{proof}[Proof of Proposition~\ref{prop:agg}]
A linear combination $\sum_i c_i\hat\theta_i$ of independent unbiased sources is unbiased iff $\sum_i c_i=1$, and its variance is $\sum_i c_i^2\sigma_i^2$. By Cauchy--Schwarz,
\[
1=\Bigl(\sum_i c_i\Bigr)^{\!2}
=\Bigl(\sum_i (c_i\sigma_i)\cdot\sigma_i^{-1}\Bigr)^{\!2}
\le\Bigl(\sum_i c_i^2\sigma_i^2\Bigr)\Bigl(\sum_i \sigma_i^{-2}\Bigr),
\]
so $\sum_i c_i^2\sigma_i^2\ge\bigl(\sum_i\sigma_i^{-2}\bigr)^{-1}=\hat\sigma_{\mathrm{agg}}^2$, with equality iff $c_i\sigma_i\propto\sigma_i^{-1}$, i.e.\ $c_i\propto\sigma_i^{-2}$---the weights in \eqref{eq:agg}. The bias expression is linearity of expectation: $\E\sum_i c_i\hat\theta_i-\theta=\sum_i c_i b_i$ with $c_i=\hat\sigma_{\mathrm{agg}}^2/\hat\sigma_i^2$.

For the contamination statement, write $\hat\theta_i=\theta+e_i+Z_i\Delta_i$, where $e_i$ is the honest (mean-zero) error, $Z_i\in\{0,1\}$ indicates contamination ($\E Z_i\le\varepsilon_i$), and $\Delta_i$ is the contaminating displacement, which by hypothesis satisfies $|\Delta_i|\le B_i$ (the contaminating support lies in a set of diameter $B_i$ containing the truth). Then
\[
\bigl|\E\hat\theta_{\mathrm{agg}}-\theta\bigr|
=\Bigl|\sum_i c_i\,\E[Z_i\Delta_i]\Bigr|
\le\sum_i c_i\,\varepsilon_i B_i
=\sum_i \varepsilon_i B_i\,\frac{\hat\sigma_{\mathrm{agg}}^2}{\hat\sigma_i^2},
\]
by the triangle inequality; no independence between the $Z_i$ is needed. The bound controls the \emph{expected} bias---the realised displacement on a contamination event can be as large as $c_iB_i$.
\end{proof}

\begin{proof}[Proof of Theorem~\ref{thm:contradiction}]
Each axis radius is nonnegative, so $\rho^{(i)}\ge 0$ always.

($\Leftarrow$) If $q^{(i)}\in\mathcal F(\hat x)$, then on every axis the unmoved value $A'=\hat A$ is feasible ($\hat x[A\mapsto\hat A]=\hat x$, since the identities already hold at the measured values), so each $\rho_A^{(i)}=0$ and $\rho^{(i)}=0$.

($\Rightarrow$) Suppose $\rho^{(i)}=0$. By \eqref{eq:rho} some axis $A$ has $\rho_A^{(i)}=0$; take a sequence $A'_n\to\hat A$ with $q^{(i)}\in\mathcal F(\hat x[A\mapsto A'_n])$. Write $(\omega_n,w_n,a_n,f_n,M_n)$ for the coordinates of $\hat x[A\mapsto A'_n]$: only the coordinates tied to axis $A$ vary, and each converges to its measured value ($a_n=1-w_n\to 1-\hat w=\hat a$ when $A=w$; $f_n=M_n/\hat N\to\hat M/\hat N=\hat f$ when $A=M$). Each feasible point carries some $\mu_n\in[\mu_{\mathrm{lo}},\mu_{\mathrm{hi}}]$ with $q^{(i)}=1-\omega_n(w_n+\mu_n(a_n-f_n))/w_n$ and $q^{(i)}\hat D\le M_n$. Pass to a subsequence with $\mu_n\to\mu^\ast$ in the compact interval; since $w_n\to\hat w>0$, eventually $w_n\ge\hat w/2>0$, so the display is continuous along the sequence and in the limit $q^{(i)}=1-\hat\omega(\hat w+\mu^\ast(\hat a-\hat f))/\hat w$ and $q^{(i)}\hat D\le\hat M$. Both membership conditions of \eqref{eq:feasible} hold at $\hat x$, hence $q^{(i)}\in\mathcal F(\hat x)$.

The second statement is immediate from the definitions: a revision of single input $A$ to $A'$ that makes the claim feasible satisfies $|A'-\hat A|/\sigma_A\ge\rho_A^{(i)}$, because $\rho_A^{(i)}$ is the infimum over exactly such revisions, and $\rho_A^{(i)}\ge\min_B\rho_B^{(i)}=\rho^{(i)}$ by \eqref{eq:rho}. The axis attaining the minimum is, by construction, the cheapest such revision in standardised units.
\end{proof}

\begin{proof}[Proof of Proposition~\ref{prop:robust}]
\emph{(i)} Condition on the realised pattern $T$. The sources outside $T$ are honest, and by the proposition's standing hypothesis every posterior in play is supported on the identified set implied by the sources it uses; a posterior formed from the honest sources plus arbitrary reports from $T$ is therefore supported on the identified set with the $T$-sources removed. Removing sources removes constraints, so this removal set contains the face-value identified set---in particular the face-value posterior support, and with it both $L_0$ and $U_0$---and it has diameter $R_T$ by definition. Any two points of a set of diameter $R_T$ are within $R_T$ of each other, so each endpoint of a support-respecting credible interval (both its endpoints lie in the posterior support, hence in the removal set) is within $R_T$ of the corresponding endpoint of $I_0$.

\emph{(ii)} The contamination indicators are independent, so pattern $T$ realises with probability $P(T)=\prod_{i\in T}\varepsilon_i\prod_{i\notin T}(1-\varepsilon_i)$, and by part (i) the endpoint displacement on that event is at most $R_T$ ($R_\emptyset=0$). Hence the expected displacement is at most $\sum_T P(T)R_T$. Splitting off the singleton patterns and bounding $\prod_{j\ne i}(1-\varepsilon_j)\le 1$,
\[
\sum_T P(T)R_T\;\le\;\sum_i\varepsilon_i R_i+\sum_{|T|\ge 2}P(T)R_T,
\]
and $P(T)\le\varepsilon_{\max}^{|T|}$ makes the second sum $O(\varepsilon^2)$ (it vanishes identically when at most one source can be contaminated). Note that $R_T$ for $|T|\ge 2$ is \emph{not} bounded by $\sum_{i\in T}R_i$ in general (two mutually corroborating sources can each have a small removal diameter while their joint removal is much larger), which is why \eqref{eq:Iepsilon} is stated as a first-order band rather than a uniform envelope.

\emph{(iii)} Write the contaminated posterior as the mixture $\hat F_\varepsilon=(1-\pi)\hat F_0+\pi G$, where $G$ is the (sub-)posterior on contaminated components and $\pi\le\tau$ by assumption. Since $|\hat F_0(x)-G(x)|\le 1$ for all $x$, $\sup_x|\hat F_\varepsilon(x)-\hat F_0(x)|\le\pi\le\tau$. Let $u\in\{\alpha/2,1-\alpha/2\}$ and abbreviate $Q_0=\hat Q_0(u)$, $Q_\varepsilon=\hat Q_\varepsilon(u)$, quantiles taken as generalised inverses. $\hat F_0$ is continuous on the segment in question (it has a density there), so $\hat F_0(Q_0)=u$, and the sup-bound $|\hat F_\varepsilon-\hat F_0|\le\tau$ holds at left limits as well. The generalised inverse satisfies $\hat F_\varepsilon(Q_\varepsilon^-)\le u\le\hat F_\varepsilon(Q_\varepsilon)$, whence
\[
\hat F_0(Q_\varepsilon)-\tau\le\hat F_\varepsilon(Q_\varepsilon^-)\le u
\le\hat F_\varepsilon(Q_\varepsilon)\le\hat F_0(Q_\varepsilon)+\tau,
\]
i.e.\ $|\hat F_0(Q_\varepsilon)-\hat F_0(Q_0)|=|\hat F_0(Q_\varepsilon)-u|\le\tau$,
and since $\hat F_0$ has density $\ge m$ on the segment between $Q_0$ and $Q_\varepsilon$, the left side is $\ge m\,|Q_\varepsilon-Q_0|$; rearranging gives $|Q_\varepsilon-Q_0|\le m^{-1}\tau$. The weight condition is substantive, not automatic: posterior mixture weights are proportional to $P(T)$ times the component marginal likelihoods, so an adversarial component that fits the observed data \emph{better} than the clean model can carry posterior weight exceeding its prior probability $\tau$.
\end{proof}

\section{The spatial Bayesian model}\label{app:spatial}

This appendix documents the forward model, priors, likelihood, and inference behind the posterior reported in Section~\ref{sec:gaza}, Step~5. Code and inputs are in the \texttt{gaza\_sim/} directory of the companion repository, and the posterior is reproducible from a fixed random seed. We emphasise its role: it is \emph{one} model-based point within the exposure-agnostic ($\mu\ge 1$) identified set, and its low posterior $q$ reflects its priors and structural choices as much as the data. Readers who reject those choices should fall back to the bounds, which do not depend on them.

\paragraph{Geometry and population.} Gaza's five governorates are split into $C=400$ cells, 80 per governorate with equal population within each governorate (cell populations therefore differ across governorates), with $N_c$ from the PCBS 2023 governorate totals. Demographic fractions (women-and-children share $w=0.733$, adult-male share $a=1-w=0.267$) are held constant across cells at the Gaza-wide values, matching the class partition of the main text.

\paragraph{Latent parameters and priors.} Each posterior draw samples, independently and uniformly on the stated ranges:

\begin{center}
\begin{small}
\begin{tabular}{@{}llll@{}}
\toprule
Parameter & Range & Meaning \\
\midrule
$M$        & $[35\mathrm{k},60\mathrm{k}]$ & total militant stock \\
$\gamma$   & $[0,1.5]$   & spatial clustering amplification of militants \\
$\sigma$   & $[0.3,1.2]$ & SD of the cell-level latent log-density \\
$\alpha$   & $[0,8]$     & strike-targeting concentration on militant density \\
$\beta$    & $[0.6,1.6]$ & targeting non-linearity exponent \\
$\mu_M$    & $[1,2.5]$   & civilian adult-male exposure multiplier \\
$\epsilon_C$ & $[0.7,1]$ & women-and-children class exposure down-weight \\
$\bar d$   & $[1.2,2.8]$ & expected kills per strike \\
$K$        & $[29\mathrm{k},40\mathrm{k}]$ & total air-to-ground strikes \\
$u$        & $[1/1.7,1]$ & recovery factor (MoH observed $/$ true direct) \\
\bottomrule
\end{tabular}
\end{small}
\end{center}

Militants are allocated across cells proportionally to $N_c\exp(\gamma u_c)$ with $u_c\sim\mathcal N(0,\sigma^2)$, capped at $30\%$ of each cell's adult-male population, with the cap and the rescaling to total $M$ iterated to a joint fixed point. (Only the product of $\gamma$ and $\sigma$ is materially identified; the two are reported jointly.)

\paragraph{Strike-kill model.} Strikes are allocated across cells with weight $N_c(1+\alpha m_c/N_c)^{\beta}$, where $m_c$ is the militant count in cell $c$. Within a cell, a death is a militant, a civilian adult man, or a woman/child with probabilities proportional to $m_c$, $(N_c a - m_c)\mu_M$, and $N_c w\,\epsilon_C$ respectively. Expected class-specific deaths are summed over cells; observed deaths are true direct deaths times the recovery factor $u$.

\paragraph{Likelihood.} Three observables: (i) the MoH cumulative total against the model's \emph{recorded} deaths $u\cdot D_{\mathrm{total}}$ (log-Gaussian, $7\%$ scale); (ii) the adult-male share among true direct deaths, $(D_{\mathrm{civAM}}+D_{\mathrm{milt}})/D_{\mathrm{total}}$, against the OHCHR verified sample (Beta likelihood at its sample size, $n=8{,}119$); and (iii) the same share against the MoH full record, down-weighted to an effective sample size of $n/5$ to reflect cruder demographic categorisation. Recording is class-neutral in the model---the recovery factor $u$ scales all classes alike---so the recorded adult-male share estimates the true share and $u$ cancels in the demographic comparison. The two demographic log-likelihoods enter with weight $1/2$ each.

\paragraph{Inference.} Importance sampling with $2\times 10^5$ prior draws (effective sample size ${\approx}3{,}500$ under the normalised weights); all posterior quantiles are weighted quantiles. The estimand is $q=D_{\mathrm{milt}}/D_{\mathrm{total}}$, the combatant share of true direct deaths, which under class-neutral recording equals the combatant share of recorded deaths---the same estimand as the main text's $q=D_M/D$. The reported posterior is $q=\nSpatialQ{}\%$ (95\% credible interval $[\nSpatialQLo{}\%,\nSpatialQHi{}\%]$), a civilian-to-combatant ratio of $\nSpatialCcr{}{:}1$ $[\nSpatialCcrLo{},\nSpatialCcrHi{}]$; the ratio is $(1-q)/q$ draw-by-draw, so the two summaries are mutually consistent by construction.

\paragraph{Sensitivities and limitations.} This model is a structured plausibility check, not independent evidence, and four caveats bound what it can show. \emph{(i) The posterior is largely prior-determined.} Under the prior alone the implied $q$ is $2.1\%$ $[1.4\%,3.0\%]$, against a posterior of $\nSpatialQ{}\%$ $[\nSpatialQLo{}\%,\nSpatialQHi{}\%]$: the likelihood barely moves $q$, because the generative structure already confines it---per person, militants carry unit kill weight, civilian adult men carry $\mu_M\in[1,2.5]$ times that weight, and the women-and-children class carries $\epsilon_C\in[0.7,1]$ (so civilian men can be \emph{more} exposed per capita than militants, and the effective male-to-W/C exposure ratio can reach ${\approx}3.6$); with no militant lethality premium, the militant share of deaths is essentially the militants' small share of the strike-weighted population. The posterior is therefore a consistency statement---the data do not push against the structure---not a data-driven estimate of $q$. \emph{(ii) Demographic blending.} The likelihood averages the OHCHR sample (at its sample size, $n=8{,}119$) and the MoH record (effective $n/5$) with weight $1/2$ each---a blending Section~\ref{sec:setup} disallows for the main analysis because the OHCHR sample over-represents residential-strike deaths \citep{spagatcockerill2024}. Re-anchoring the demographic likelihood entirely on the MoH record moves the posterior to $2.0\%$ $[1.4\%,2.8\%]$; entirely on OHCHR, to $2.3\%$ $[1.6\%,3.1\%]$: either anchor stays far inside the identified set. \emph{(iii) The spatial structure is latent.} No observed coordinates, strike locations, or cell-level casualty counts enter the likelihood, which sees only the aggregate toll and the two demographic shares; ``spatial'' describes the generative model, not fitted spatial evidence. \emph{(iv) Window mismatch.} The strike-count prior is sourced through January 2025 while the toll target is late-2025, so the per-strike kill parameter absorbs the gap. Unlike the exposure calibration of Section~\ref{sec:bounds}, the model has not been validated on historical conflicts with known combatant status. Because of all four caveats, the paper's claims do not rest on this posterior; it appears as one row of Table~\ref{tab:convergence}, flagged as prior-dependent.

\section{The 81-conflict dataset}\label{app:conflicts}

Table~\ref{tab:all-conflicts} lists every conflict in the dataset behind the overview of Section~\ref{sec:empirical}, with midpoint death totals and the implied civilian-share intervals. The table is generated by the same script that draws Figures~\ref{fig:range-frame} and~\ref{fig:uncertainty-ladder}. The intervals are those of the curated per-conflict sources under each conflict's own conventions; in particular, the Gaza row carries the dataset's central range (informed by the spatial posterior of Appendix~\ref{app:spatial}), which is narrower than the identified-set bounds of Section~\ref{sec:gaza}---the bounds are the paper's claim, the row is the dataset's summary. The supplementary material provides a single master table detailing the bounds, binding relations, and confidence grades for all 81 conflicts, along with a complete list of primary sources used for each.

% Auto-generated by paper/build_figures.py -- do not edit by hand.
\begin{footnotesize}
\renewcommand{\arraystretch}{1.2}
\begin{longtable}{@{}lccrc@{}}
\caption{\textbf{The 81-conflict dataset behind Figures~\ref{fig:range-frame} and~\ref{fig:uncertainty-ladder}} (conflicts with at least 1{,}000 attributed deaths shown). Deaths are midpoints of the curated total-death ranges (military $+$ civilian, including indirect where sources attribute it); the civilian-share interval spans the most and least civilian-heavy readings of the source ranges. Where sources state no ceiling for one casualty class, the ceiling is closed mechanically by the grand-total residual (total high minus the other class's floor) rather than treated as zero. Per-conflict sources are in the supplement.}\label{tab:all-conflicts}\\
\toprule
Conflict & Period & Deaths (mid) & Civilian share & lo--hi \\
\midrule
\endfirsthead
\multicolumn{5}{@{}l}{\footnotesize\emph{Table~\ref{tab:all-conflicts} continued}}\\
\toprule
Conflict & Period & Deaths (mid) & Civilian share & lo--hi \\
\midrule
\endhead
\bottomrule
\endfoot
Boxer Rebellion & 1899--1901 & 125k & 89\% & 63--98\% \\
Philippine-American War & 1899--1902 & 211k & 89\% & 86--92\% \\
Herero and Namaqua Genocide & 1904--1908 & 71.6k & 98\% & 97--98\% \\
Russo-Japanese War & 1904--1905 & 167k & 12\% & 11--14\% \\
Mexican Revolution & 1910--1920 & 1.25M & 84\% & 50--93\% \\
First and Second Balkan Wars & 1912--1913 & 1.31M & 82\% & 72--86\% \\
Armenian Genocide and late Ottoman genocide\dots{} & 1914--1923 & 1.88M & 100\% & 100--100\% \\
World War I & 1914--1918 & 18.5M & 50\% & 37--59\% \\
Russian Civil War & 1917--1923 & 9.5M & 88\% & 77--94\% \\
Greco-Turkish War and population exchange & 1919--1922 & 885k & 94\% & 81--97\% \\
Irish War of Independence and Civil War & 1919--1923 & 4.22k & 39\% & 36--42\% \\
Rif War & 1921--1927 & 84k & 22\% & 13--30\% \\
Chaco War & 1932--1935 & 92.5k & 12\% & 1--24\% \\
Soviet repression: Holodomor and Great Purge & 1932--1939 & 8.44M & 100\% & 100--100\% \\
The Holocaust & 1933--1945 & 14.1M & 77\% & 71--81\% \\
Second Italo-Ethiopian War & 1935--1937 & 609k & 62\% & 49--84\% \\
Spanish Civil War & 1936--1939 & 402k & 56\% & 43--67\% \\
Second Sino-Japanese War & 1937--1945 & 22.6M & 84\% & 68--92\% \\
World War II - European/African/Atlantic Th\dots{} & 1939--1945 & 42M & 56\% & 51--61\% \\
World War II - Pacific Theater/Asia-Pacific\dots{} & 1941--1945 & 25M & 76\% & 69--82\% \\
Bengal Famine of 1943 & 1943--1944 & 2.65M & 100\% & 100--100\% \\
Chinese Civil War & 1945--1949 & 6.01M & 83\% & 81--85\% \\
Indonesian National Revolution & 1945--1949 & 164k & 48\% & 19--72\% \\
First Indochina War & 1946--1954 & 621k & 43\% & 22--60\% \\
1948 Arab-Israeli War/Nakba & 1947--1949 & 24.2k & 52\% & 42--62\% \\
Partition of India & 1947--1948 & 1.1M & 100\% & 95--100\% \\
Malayan Emergency & 1948--1960 & 11.4k & 25\% & 22--28\% \\
Korean War & 1950--1953 & 3.24M & 71\% & 58--80\% \\
Mau Mau Uprising & 1952--1960 & 37.5k & 54\% & 29--77\% \\
Algerian War of Independence & 1954--1962 & 510k & 66\% & 48--75\% \\
Vietnam War & 1955--1975 & 2.78M & 55\% & 32--68\% \\
Great Leap Forward famine & 1958--1962 & 30M & 100\% & 100--100\% \\
Congo Crisis & 1960--1965 & 182k & 67\% & 36--98\% \\
Guatemalan Civil War & 1960--1996 & 195k & 98\% & 95--100\% \\
North Yemen Civil War & 1962--1970 & 150k & 38\% & 0--67\% \\
Indonesian mass killings of 1965-66 & 1965--1966 & 600k & 100\% & 100--100\% \\
Nigerian Civil War & 1967--1970 & 1.77M & 96\% & 83--98\% \\
Six-Day War & 1967--1967 & 17k & 7\% & 0--17\% \\
Bangladesh Liberation War/1971 genocide & 1971--1971 & 2.18M & 99\% & 89--100\% \\
Yom Kippur War & 1973--1973 & 16k & 1\% & 0--2\% \\
Ethiopian Civil War/Derg/1983-85 famine & 1974--1991 & 1.35M & 49\% & 30--89\% \\
Angolan Civil War & 1975--2002 & 650k & 84\% & 68--95\% \\
Indonesian occupation of East Timor & 1975--1999 & 154k & 96\% & 90--98\% \\
Khmer Rouge Cambodian Genocide & 1975--1979 & 2.03M & 99\% & 98--99\% \\
Lebanese Civil War & 1975--1990 & 115k & 62\% & 41--93\% \\
Mozambican Civil War & 1977--1992 & 800k & 95\% & 92--97\% \\
Cambodian-Vietnamese War & 1978--1989 & 213k & 72\% & 65--80\% \\
Salvadoran Civil War & 1979--1992 & 75k & 65\% & 60--70\% \\
Soviet-Afghan War & 1979--1989 & 1.06M & 89\% & 86--91\% \\
Iran-Iraq War & 1980--1988 & 524k & 21\% & 14--29\% \\
Second Sudanese Civil War & 1983--2005 & 2.06M & 98\% & 97--98\% \\
Sri Lankan Civil War & 1983--2009 & 90k & 42\% & 28--53\% \\
First Intifada & 1987--1993 & 2.21k & 97\% & 97--97\% \\
Nagorno-Karabakh wars & 1988--2023 & 39k & 28\% & 4--50\% \\
Gulf War & 1990--1991 & 208k & 66\% & 34--91\% \\
Rwandan Genocide and Rwandan Civil War & 1990--1994 & 830k & 99\% & 99--99\% \\
Somali Civil War & 1991-- & 675k & 90\% & 71--97\% \\
Yugoslav Wars - Bosnia/Croatia & 1991--1995 & 123k & 40\% & 38--42\% \\
First Chechen War & 1994--1996 & 56.2k & 71\% & 56--86\% \\
First Congo War & 1996--1997 & 173k & 60\% & 20--100\% \\
Kosovo War & 1998--1999 & 13k & 76\% & 74--77\% \\
Second Congo War & 1998--2003 & 3.65M & 100\% & 93--100\% \\
Second Chechen War & 1999--2009 & 56.2k & 67\% & 48--82\% \\
Second Intifada & 2000--2005 & 4.95k & 65\% & 53--75\% \\
War in Afghanistan & 2001--2021 & 268k & 41\% & 20--57\% \\
Darfur conflict/genocide & 2003--2020 & 235k & 78\% & 57--100\% \\
Iraq War & 2003--2011 & 457k & 90\% & 72--94\% \\
Mexican Drug War & 2006-- & 345k & 62\% & 29--87\% \\
Boko Haram insurgency/Lake Chad conflict & 2009-- & 375k & 89\% & 83--94\% \\
Libyan Civil Wars & 2011--2020 & 31.5k & 50\% & 13--75\% \\
Syrian Civil War & 2011--2024 & 618k & 44\% & 40--47\% \\
South Sudanese Civil War & 2013--2020 & 380k & 85\% & 75--95\% \\
War against ISIS in Iraq and Syria & 2014--2019 & 248k & 59\% & 53--67\% \\
War in Donbas & 2014--2022 & 14.3k & 24\% & 24--24\% \\
Yemeni Civil War & 2014-- & 414k & 65\% & 53--78\% \\
Rohingya genocide/Myanmar military operatio\dots{} & 2016-- & 27.3k & 98\% & 96--99\% \\
Tigray War & 2020--2022 & 431k & 57\% & 35--79\% \\
Myanmar civil war & 2021-- & 104k & 39\% & 17--60\% \\
Russian invasion of Ukraine & 2022-- & 521k & 11\% & 3--26\% \\
Israel-Hamas/Gaza war & 2023-- & 96.4k & 96\% & 94--98\% \\
Sudan war & 2023-- & 230k & 93\% & 49--97\% \\
\end{longtable}
\end{footnotesize}

